% \input{Dissertation/appendixsetup}   % Предварительные настройки для правильного подключения Приложений

\chapter{Фрагменты кода встроенного ПО контроллера верхнего уровня} \label{AppendixA}

В приведенном ниже листинге \ref{list:restrict} представлен код алгоритма ограничения координат УСП.

\lstinputlisting[language={[ISO]C++}, caption={Листинг кода ограничения координат УСП},label={list:restrict}]{listings/gold_section_algorithm.cpp}

%В листинге \ref{list:inverseKin} приведен код решения обратной задачи для 3DOF динамической платформы
%\lstinputlisting[language={[ISO]C++}, caption={Листинг кода решения ОЗК},label={list:inverseKin}]{listings/Solver.cpp}
%
%В листинге ниже приведен код решения ПЗД для 3DOF динамической платформы.
%\lstinputlisting[language={[ISO]C++}, caption={Листинг кода решения ПЗД},label={list:solverDyn}]{listings/dynamicsolver.cpp}

\chapter{Характеристики планшетного компьютера АТМ12} \label{AppendixB}
Контроллер движения выполнен на базе специализированного планшетного ПК АТМ12, показанного на рисунке \ref{img:ATM12pic}. Компьютер АТМ12 предназначен для использования в системах АСУ в качестве устройства управления и отображения данных.

\begin{figure}
	\center
	\includegraphics[width=0.5\linewidth]{ATM12pic}
	\caption{Контроллер движения АТМ12}
	\label{img:ATM12pic}
\end{figure}

\begin{table} [htbp]%
    \centering
	\caption{Характеристики аппаратной части контроллера АТМ12}%
	\label{tbl:ATM12parameters}
    \renewcommand{\arraystretch}{1.5}
    \begin{SingleSpace}
	\begin{tabulary}{\textwidth}{@{}>{\zz}L >{\zz}L >{\zz}L >{\zz}L >{\zz}L@{}}
        \toprule     %%% верхняя линейка
    	п/п	& Параметр & Значение для АТМ12\\
        \midrule
    	1 & Платформа & ARM Cortex-A9 Nvidia Tegra 2 (1.4GHz 1MB Cache) \\
    	2 & Операционная система & Linux (RT)\\
    	3 & ОЗУ & 4GB DDR2\\
    	4 & ПЗУ & 2GB NAND Flash; microSDH до 64GB x 2\\
    	5 & Экран и графика & 12.1'' (TFT), 768x1024, сенсорный резистивный экран\\
    	6 & Интерфейсы & LAN: x1, 4pin M12; CAN/DC: x1 5pin M12; USB2.0 защищенный x2\\
    	7 & Питание & Блок питания переменного тока 100–240V AC либо внешний источник постоянного тока 12V(9-36V) DC. Потребляемая мощность: 12 Вт в нормальном режиме \\
    	8 & Защита информации & Программная защита: ограничение по паролю (администратор, пользователь, и доступ к ПЗУ)  Защита от НСД собственной разработки \\
    	9 & Защита от внешних воздействий & Защита от воды:IPX5  Защита от пыли:IP7X  Защита при вибрации и тряски: 5Hz-500Hz/1Grms/3Axis  Рабочая температура: -40...+70 $^oC$  Рабочая относительная влажность 5\%--95\%\\
    	10 & Охлаждение & Пассивное герметичное (на корпус) \\
        \bottomrule
	\end{tabulary}%
   	\end{SingleSpace}
\end{table}

АТМ12 работает под управлением операционный системы Linux с открытым исходным кодом, что позволяет пройти сертификацию устройства на отсутствие программных <<закладок>>. Сборка операционной системы и разработанное ПО:
\begin{itemize}
	\item может использоваться длительное время без изменений,
	\item адаптируется под конкретную задачу управления СП,
	\item исключает использование устройства не по назначению,
	\item способствует более надежной защите устройства от несанкционированного доступа.
\end{itemize}
Наличие промышленных интерфейсов и программная поддержка распространенных протоколов (Modbus, CANOpen) позволяет использовать АТМ12 для широкого круга задач управления подвижными объектами.

\chapter{Акт о внедрении результатов диссертационного исследования}
\begin{figure}
\center
\includegraphics[width=0.8\linewidth]{actRealization1}
\end{figure}

\begin{figure}
\center
\includegraphics[width=\linewidth]{actRealization2}
\end{figure}